\documentclass[12pt,a4paper]{article}

\usepackage{amsmath}
\usepackage{amssymb}
\usepackage[margin=1in]{geometry}
\usepackage{CJKutf8}

\begin{document}
\begin{CJK*}{UTF8}{bsmi}

\noindent \textbf{試題 :}

\begin{itemize}
    \item Please give details of your calculation. A direct answer without explanation is not counted.
    \item Your answers must be in English.
\end{itemize}

\vspace{0.5cm}

\section*{Problem 1 (15\%)}
What is the bit string of $-37$ under IEEE 754 standard? We assume single precision. Please explain details instead of just giving a string.

\vspace{0.5cm}

\section*{Problem 2 (30\%)}
Consider the following matrix:
\[
A = \begin{bmatrix}
1 & 4 & 7 \\
2 & 5 & 8 \\
3 & 6 & 10
\end{bmatrix}
\]
We would like to find the LU factorization using the pivoting.

\begin{enumerate}
    \item[(a)] Give all details of finding
    \[ M_{n-1} P_{n-1} \dots M_1 P_1 A = U \]
    
    \item[(b)] Find $P$ such that
    \begin{equation}
    PA = LU \label{eq1}
    \end{equation}
    where $L$ and $U$ are lower- and upper-triangular matrices, respectively.
    
    \item[(c)] If we solve the linear system
    \[
    Ax = b = \begin{bmatrix} 6 \\ 9 \\ 12 \end{bmatrix}
    \]
    show how to use (1) to solve the linear system.
\end{enumerate}

\vspace{0.5cm}

\section*{Problem 3 (20\%)}
Consider the following matrix:
\[
A = \begin{bmatrix}
16 & 8 & 4 & 12 \\
8 & 8 & 0 & 2 \\
4 & 0 & 27 & 0 \\
12 & 2 & 0 & 30
\end{bmatrix}
\]
Generate the Cholesky factorization using the outer product form.

\vspace{0.5cm}

\section*{Problem 4 (20\%)}
In exact arithmetic, the addition operation is commutative, i.e.,
\[ x + y = y + x \]
for any two numbers $x, y$ and also associative, i.e.,
\[ x + (y + z) = (x + y) + z \]
for any $x, y$ and $z$. Is the floating point addition operation commutative? Is it associative? We assume exact rounded operations.

\vspace{0.5cm}

\section*{Problem 5 (15\%)}
What is the smallest positive denormalized number under the IEEE standard (single precision)?

\end{CJK*}
\end{document}